\section*{Devoir maison }

On s'intéresse au code suivant :

\begin{lstlisting}[language=Python]
class Pile:
    ''' 
    On cherche à implémenter l'interface de la Pile sous forme de classe
    '''
    def __init__(self, tab):
        self.pile = tab

    def est_vide(self):
        return len(self.pile) == 0

    def empiler(self, x):
        ......................

    def depiler(self):
        return .................

    def tete(self):
        return self.pile[-1]
                       
\end{lstlisting}

\begin{Exercise}[title={POO - pile}]
    \Question Quels sont les attributs de cette classe ?
    \Question Quelles sont les méthodes de cette classe ?
    \Question De quelle ligne à quelle ligne se trouve le constructeur ?
    \Question Donner l'instruction permettant de créer une variable $p1$ qui est une pile vide.\\
    \emptybox{1cm}
    \Question Donner l'instruction permettant de créer une variable $p2$ qui est la pile [5, 8, 6, 5, 3] avec 3 comme tête de pile.\\
    \emptybox{1cm}
    \Question Que vaut p2.depiler() + p2.depiler() + p2.tete() ?\\
    \emptybox{1cm}
    \Question Compléter la méthode empiler.
    \Question Compléter la méthode depiler.
    \Question  Donner le code permettant de :
    \begin{itemize}
        \item Créer une pile p vide :
        \item Empiler 2 dans p :
        \item Tester si p est vide :
        \item Dépiler p :
        \item Empiler dans p une copie de l'élément en tête de p :
    \end{itemize}
\end{Exercise}